\begin{explanationbox}
	\textbf{\textcolor{CExplanation}{一句话直觉}}

	椭圆周长没有简单的封闭表达式，但拉马努金公式用一个巧妙的修正因子给出了极高精度的近似结果，相当于“圆周长乘以一个形状修正项”。

	\medskip

	\textbf{\textcolor{CExplanation}{核心拆解}}

	\begin{description}[style=nextline, leftmargin=2.8em, labelsep=0.5em, itemsep=0.45em]
		\item[\textbf{要点一（h参数含义）：}]
		      $h=\left(\frac{a-b}{a+b}\right)^2$ 刻画椭圆扁平程度，$h=0$ 对应圆，$h$ 越大椭圆越扁。

		\item[\textbf{要点二（公式结构）：}]
		      公式形如 $C\approx \pi(a+b)\times\left(1+\frac{3h}{10+\sqrt{4-3h}}\right)$，整体为“半平均周长×修正因子”。

		\item[\textbf{要点三（根号作用）：}]
		      分母中的 $\sqrt{4-3h}$ 项保证当 $h$ 较大时分母不过小，从而维持近似稳定，避免高离心率时发散。
	\end{description}

	\medskip

	\textbf{\textcolor{CExplanation}{几何本质（椭圆形状修正）}}

	椭圆周长在几何上无法用初等函数表示，该公式通过将椭圆视为圆在长轴方向拉伸后，再对周长进行非线性修正，本质上是一个关于偏心率 $e$ 的有理逼近。

	\medskip

	\textbf{\textcolor{CExplanation}{代数意义（连分数逼近）}}

	该公式可以看作第二类完全椭圆积分 $E(e)$ 的一个代数逼近，通过连分数展开截断得到，具有紧凑且高精度的形式。

	\medskip

	\textbf{\textcolor{CExplanation}{考点价值（高考应用）}}

	\begin{itemize}[leftmargin=*, itemsep=0.5em, topsep=0.4em]
		\item \textbf{考法一（直接代入）：} 已知 $a,b$ 直接代入公式算近似周长，需仔细计算 $h$ 与根号。
		\item \textbf{考法二（参数反推）：} 已知 $C$ 与 $a$（或 $b$），可反解另一参数的近似值，通常转化为方程求解。
	\end{itemize}

	\medskip

	\textbf{\textcolor{CExplanation}{顿悟点}}

	\textbf{当 $a$ 与 $b$ 相差不大时（如 $a/b<2$），该公式误差极小（小于 $0.04\%$），完全可以作为精确公式使用，极大简化计算。}

	\medskip

	\textbf{\textcolor{CExplanation}{使用场景}}

	\begin{itemize}[leftmargin=*, itemsep=0.5em, topsep=0.4em]
		\item \textbf{场景一（快速估算）：} 已知椭圆方程，快速计算周长用于比较或实际估算。
		\item \textbf{场景二（题干给定）：} 在考试中若题干给出该公式，直接应用，无需重新推导。
	\end{itemize}
\end{explanationbox}